% ====================================================================
% CAPÍTULO 4: RESULTADOS
% ====================================================================
\chapter{Resultados}
% Nota de avance: En este capítulo se presentan los datos cuantitativos teóricos.
% obtenidos probablemente a través del modelamiento y la distribución espacial.

% Los resultados del modelamiento electroquímico confirman la viabilidad de la hidrólisis a profundidades abisales. Como se observa en la Tabla \ref{tab:datos_oxigeno}, existe una correlación directa entre la densidad estimada de los nódulos polimetálicos y el flujo de oxígeno oscuro generado en el lecho oceánico.

% --- Inserción de Tabla  ---
%\begin{table}[h]
\centering
\small % (achica un poco la letra)
\begin{tabular}{|l|c|c|}
\hline
\textbf{Zona Abisal} & \textbf{Densidad de Nódulos ($kg/m^2$)} & \textbf{Flujo de $O_2$ ($mmol/m^2/d$)} \\ \hline
Clarion-Clipperton & 15.4 & 0.85 \\ \hline
Cuenca del Perú & 8.2 & 0.41 \\ \hline
Pacífico Sur & 5.1 & 0.22 \\ \hline
\end{tabular}
\caption{Estimación teórica del flujo de oxígeno oscuro según densidad nodular regional.}
\label{tab:datos_oxigeno}
%\end{table}

% Por otro lado, la distribución espacial de las zonas de mayor potencial de óxido-reducción coincide fuertemente con las anomalías geomagnéticas del fondo marino. En la Figura \ref{fig:mapa_nodos} se presenta un esquema referencial de estas áreas de alta concentración, listas para ser mapeadas en detalle mediante sistemas de información geográfica.

% --- Inserción de Figura  ---
%\begin{figure}[h]
\centering
% Nota: Añadimos otra imagen en la carpeta de imagenes para esta parte
\includegraphics[width=0.4\textwidth]{imagenes/}
\caption{Mapa referencial de anomalía geomagnética global.}
\label{fig:mapa_nodos}
%\end{figure}